%%%%%%%%%%  scan idx 159  |  folio 146  %%%%%%%%%%
\SGAPrintPage{146}%
(the two on the left being homomorphisms of presheaves of rings), which, by (6.7), induce isomorphisms on the associated sheaves. We shall show that the associated sheaves are in fact canonically isomorphic to the constant sheaf \(\mathbb Z_{\underline{S}}\).

\sgaStableTarget{sga6.e01.lem.6-10}\textbf{Lemma 6.10.} \uline{Let} \((\underline{S},A)\) \uline{be a locally ringed topos,} \(U\) \uline{an object of} \(\underline{S}\) \uline{distinct from the empty object, and} \(p,q\in\mathbb N\)\uline{. If we have} \(A_U^p\simeq A_U^q\)\uline{, then} \(p=q\)\uline{.}

\uline{Proof}. By an argument similar to the one given in the proof of (2.15.1), we reduce to the case where \((\underline{S},A)\) is the locally ringed topos defined by a prescheme, and in this case the assertion is trivial (take a point in \(U\)).

\sgaStableTarget{sga6.e01.def.6-11}\textbf{Definition 6.11.} \uline{We call} rank, \uline{and denote by} \(\rg\), \uline{the additive} \(\underline{S}\)\uline{-function from} \(\underline{C}_1\) \uline{to} \(\mathbb Z_{\underline{S}}\) \uline{defined by}
\[
  \rg(L)=p\quad\text{if } L\simeq A_X^p,
  \qquad X\ne\varnothing,
\]
\[
  \rg(L)=0\quad\text{if } L\in\Ob\underline{C}_{1X},
  \qquad X=\varnothing.
\]
By (6.7), the rank function extends uniquely to an additive \(\underline{S}\)-function from \(\Parf(\underline{S})\) to \(\mathbb Z_{\underline{S}}\) (which we shall again call rank and denote by \(\rg\)).

\sgaStableTarget{sga6.e01.prop.6-12}\textbf{Proposition 6.12.} a) \uline{For} \(L,M\in\Ob\Dpx^-(X)_{\mathrm{parf}}\) (\(X\in\Ob \underline{S}\)), \uline{we have} \(\rg(L\overset{\mathrm L}{\otimes} M)=\rg(L)\,\rg(M)\).

b) \uline{Let} \(u:(\underline{S}',A')\longrightarrow(\underline{S},A)\) \uline{be a morphism of locally ringed topoi. For} \(M\in\Ob\Dpx^-(\underline{S})_{\mathrm{parf}}\), we have \(\rg(\mathbb L u^*(M))=u^*(\rg(M))\). \uline{In other words, the rank homomorphism}: \(k(\Dpx^-(\underline{S})_{\mathrm{parf}})\longrightarrow \mathbb Z_{\underline{S}}\) \uline{is a homomorphism of presheaves of rings, compatible with inverse images}.

% Additional notation for Exposé I
\providecommand{\Cpx}{\mathrm C}
\providecommand{\Kpx}{\mathrm K}
\providecommand{\Dpx}{\mathrm D}
\providecommand{\Ho}{\mathrm H}
\providecommand{\RHom}{\mathrm R\!\operatorname{Hom}}
\renewcommand{\Ob}{\operatorname{ob}}
\providecommand{\Hom}{\operatorname{Hom}}
\providecommand{\RHom}{\mathrm R\!\operatorname{Hom}}
\providecommand{\Ext}{\operatorname{Ext}}
\providecommand{\Tor}{\operatorname{Tor}}
\providecommand{\Tr}{\operatorname{Tr}}
\providecommand{\Add}{\operatorname{Add}}
\providecommand{\Cart}{\underline{\operatorname{Cart}}}
\providecommand{\Parf}{\operatorname{Parf}}
\providecommand{\coh}{\operatorname{coh}}
\providecommand{\amp}{\operatorname{amp}}
\providecommand{\cl}{\operatorname{cl}}
\providecommand{\rg}{\operatorname{rg}}
\providecommand{\Id}{\operatorname{Id}}
\providecommand{\Ker}{\operatorname{Ker}}
\renewcommand{\Im}{\operatorname{Im}}
\providecommand{\cH}{\mathcal H}
\providecommand{\Zs}{\mathbb Z_S}
\providecommand{\A}{\mathcal A}
\providecommand{\cC}{\mathcal C}
\providecommand{\cCzero}{\mathcal C_0}
\providecommand{\otL}{\mathop{\otimes}^{\mathrm L}}
\providecommand{\tamp}{\operatorname{tor.amp}}
\providecommand{\pamp}{\operatorname{parf.amp}}
\providecommand{\pdim}{\operatorname{parf.dim}}
\providecommand{\bon}{\mathrm b}
\providecommand{\parf}{\operatorname{parf}}
\providecommand{\parff}{\operatorname{parf}}
\providecommand{\ac}{\mathrm{ac}}


% Additional notation for Exposé I, completion pages 154--166
\providecommand{\End}{\operatorname{End}}
\providecommand{\tr}{\operatorname{tr}}

%%%%%%%%%%  scan idx 160  |  folio 147  %%%%%%%%%%
\SGAPrintPage{147}%
\uline{Proof}. It suffices to show that the homomorphism induced by rank on the sheaf associated to \(k(\Dpx^-(\underline{S})_{\mathrm{parf}})\) is a homomorphism of sheaves of rings compatible with inverse images. By (6.9), we can replace \(k(\Dpx^-(\underline{S})_{\mathrm{parf}})\) by \(k(\underline{C}_1)\). We are therefore reduced to checking a) and b) for \(L\) and \(M\) free of finite type, which is trivial.

\sgaStableTarget{sga6.e01.para.6-13}6.13. With \((\underline{S},A)\) still denoting a locally ringed topos, let \(X\in\Ob \underline{S}\) and \(\underline n\in H^0(X,\mathbb Z_{\underline{S}})\). We define an object of \(\Parf(X)\), denoted by \(A_X^{\underline n}\), as follows. When \(\underline n\) is constant with value an integer \(n\geq0\) (resp.\ \(n<0\)), we put \(A_X^{\underline n}=A_X^n\) (resp.\ \(A_X^{\underline n}=A_X^{-n}[1]\)); in the general case, we decompose \(X\) into disjoint pieces \(X_i\) such that \(\underline n|X_i\) is constant, and \(A_X^{\underline n}\) is defined as the unique object of \(\Parf(X)\) inducing \(A_{X_i}^{\underline n|X_i}\) on each \(X_i\). (More generally, we can define \(L^{\underline n}\) for \(L\in\Ob\Parf(X)\), by putting \(L^{\underline n}=A_X^{\underline n}\otimes L\).).
The map \(i_X:H^0(X,\mathbb Z_{\underline{S}})\longrightarrow k(\Parf(X))\), defined by \(i_X(\underline n)=\cl(A_X^{\underline n})\), is a ring homomorphism (such that \(i_X(1)=1\)), and defines, as \(X\) varies, a homomorphism of presheaves of rings
\[
  i:\mathbb Z_{\underline{S}}\longrightarrow k(\Parf(\underline{S})).
\]

\sgaStableTarget{sga6.e01.prop.6-14}\textbf{Proposition 6.14.} \uline{Under the hypotheses of} (6.13), \uline{we have}
\[
  \rg\, i=\Id(\mathbb Z_{\underline{S}})
\]
(\uline{in other words,} \(k(\Parf(\underline{S}))\)\uline{, considered as a presheaf of} \(\mathbb Z_{\underline{S}}\)\uline{-algebras by means of} \(i\)\uline{, is augmented toward} \(\mathbb Z_{\underline{S}}\) \uline{by rank}). \uline{Moreover, rank}
\[
  \rg:k(\Parf(\underline{S}))\longrightarrow\mathbb Z_{\underline{S}}
\]
\uline{defines an isomorphism on the associated sheaves}.

%%%%%%%%%%  scan idx 161  |  folio 148  %%%%%%%%%%
\SGAPrintPage{148}%
\uline{Proof}. The first assertion follows trivially from the relations \(\cl(L[1])=-\cl(L)\), \(\cl(L^n)=n\cl(L)\) for \(L\in\Ob\Parf(\underline{S})\). Let us prove the second. First observe that \(i\) factors in the evident way through \(k(\underline{C}_0)\). Since \(\rg\, i\) is the identity, we are reduced (by (6.9)) to proving that the homomorphism \(i\, \rg:k(\underline{C}_0)\longrightarrow k(\underline{C}_0)\) induces an isomorphism on the associated sheaves, i.e.\ that, for \(L\in\Ob\underline{C}_0\), \(i\, \rg(L)\) is locally isomorphic to \(L\), which is immediate.

We deduce, by (6.6.b)),

\sgaStableTarget{sga6.e01.cor.6-15}\textbf{Corollary 6.15.} \uline{Under the hypotheses of} (6.13), \uline{for every abelian sheaf} \(F\) \uline{on} \(\underline{S}\)\uline{, the group of additive} \(\underline{S}\)\uline{-functions from} \(\underline{\Parf}(\underline{S})\) \uline{to} \(F\) \uline{is canonically identified with} \(H^0(S,F)\) (\(S\) = \uline{final object of} \(\underline{S}\)), \uline{the rank function corresponding to} \(F=\mathbb Z_{\underline{S}}\) \uline{and to the unit section of} \(\mathbb Z_{\underline{S}}\).

\sgaStableTarget{sga6.e01.sec.7}\section*{7. Duality of perfect complexes}

\sgaStableTarget{sga6.e01.para.7-0}7.0. In this section, \((\underline{S},A)\) denotes a topos ringed in commutative rings; we denote by \(S\) the final object of \(\underline{S}\). We propose to generalize to perfect complexes the notions of duality well known for locally free sheaves of finite type.

\sgaStableTarget{sga6.e01.prop.7-1}\textbf{Proposition 7.1.} \uline{Let} \(E\in\Ob\Dpx(S)_{\mathrm{parf}}\), \(F\in\Ob\Dpx^+(S)\).

a) \uline{When} \(F\) \uline{has locally bounded cohomology, the same holds for} \(\mathbb{R}\underline{\Hom}(E,F)\). \uline{When} \(F\in\Ob\Dpx(S)\) \uline{is pseudo-coherent (resp.\ locally of finite tor-dimension, resp.\ perfect), the same holds for} \(\mathbb{R}\underline{\Hom}(E,F)\).

%%%%%%%%%%  scan idx 162  |  folio 149  %%%%%%%%%%
\SGAPrintPage{149}%
b) Suppose that \(E\) is moreover of finite tor-dimension (a condition satisfied when \(\underline{S}\) is of finite type). Then, when \(F\) has bounded cohomology (resp.\ finite tor-dimension), the same holds for \(\mathbb{R}\underline{\Hom}(E,F)\). More precisely, let \(a,b,m,n\) be integers such that \(\tamp(E)\subset[a,b]\) (or, equivalently, by (5.8), \(\pamp(E)\subset[a,b]\)) and \(\Ho^i(F)=0\) for \(i\notin[m,n]\) (resp.\ \(\tamp(F)\subset[m,n]\)). Then we have \(\Ho^i(\mathbb{R}\underline{\Hom}(E,F))=0\) for \(i\notin[m-b,n-a]\) (resp.\ \(\tamp(\mathbb{R}\underline{\Hom}(E,F))\subset[m-b,n-a]\)).

\uline{Proof}. Assertion a) is local, so we can suppose that \(E\) is strictly perfect, and we conclude by dévissage on \(E\) (reducing to \(E=A\)). For assertion b), we can suppose that \(E\) is strictly perfect with \(E^i=0\) for \(i\notin[a,b]\), and \(F^i=0\) for \(i\notin[m,n]\) (resp.\ \(F^i\) flat for every \(i\) and \(F^i=0\) for \(i\notin[m,n]\)). It then suffices to apply the following lemma:

\sgaStableTarget{sga6.e01.lem.7-1-1}\textbf{Lemma 7.1.1.} \uline{Let} \(E\in\Ob\Dpx^-(S)\), \(F\in\Ob\Dpx^+(S)\). \uline{When} \(E\) \uline{is strictly pseudo-coherent} (i.e.\ (2.1) \uline{with components direct factors of free Modules of finite type}), \uline{the canonical arrow of} \(\Dpx^+(S)\)
\[
  \underline{\Hom}^\bullet(E,F)\longrightarrow\mathbb{R}\underline{\Hom}(E,F)
\]
\uline{is an isomorphism}.

\uline{Proof}. The spectral sequence
\[
  E_2^{pq}=\underline{\Ho}^p(\underline{\Ext}^q(E,F)^\bullet)\Longrightarrow \underline{\Ext}^*(E,F),
\]
where
\[
  \underline{\Ext}^q(E,F)^i=\coprod_{i_2-i_1=i}\underline{\Ext}^q(E^{i_1},F^{i_2}),
\]
degenerates since \(\underline{\Ext}^q(E,F)^\bullet=0\) for every \(q>0\) (1.3.1). We deduce isomorphisms
%%%%%%%%%%  scan idx 163  |  folio 150  %%%%%%%%%%
\SGAPrintPage{150}%
\[
  E_2^{p0}=\underline{\Ho}^p(\underline{\Hom}^\bullet(E,F))
  \xrightarrow{\sim}\underline{\Ext}^p(E,F)=\underline{\Ho}^p(\mathbb{R}\underline{\Hom}(E,F)),
\]
which proves the assertion.

\sgaStableTarget{sga6.e01.prop.7-1-2}\textbf{Proposition 7.1.2.} \uline{Let} \(f:(\underline{S}',A')\longrightarrow(\underline{S},A)\) \uline{be a morphism of ringed topoi. Let} \(E\in\Ob\Dpx^-(S)\) \uline{and} \(F\in\Ob\Dpx^b(S)\) \uline{be such that} \(\mathbb{R}\underline{\Hom}(E,F)\in\Ob\Dpx^b(S)\) \uline{and} \(\mathbb{L}f^*(F)\in\Ob\Dpx^b(S')\). \uline{There exists a canonical functorial arrow}
\[
  \mathbb{L}f^*\mathbb{R}\underline{\Hom}(E,F)
  \longrightarrow
  \mathbb{R}\underline{\Hom}(\mathbb{L}f^*E,\mathbb{L}f^*F).
\]
\uline{It is an isomorphism when} \(E\) \uline{is perfect}.

\uline{Proof}. By the adjunction of the functors \(\mathbb{L}f^*\) and \(\mathbb{R}f_*\), it suffices to define an arrow
\[
  \mathbb{R}\underline{\Hom}(E,F)
  \longrightarrow
  \mathbb{R}f_*\mathbb{R}\underline{\Hom}(\mathbb{L}f^*E,\mathbb{L}f^*F).
\]
For this we use the adjunction arrow \(F\to\mathbb{R}f_*\mathbb{L}f^*F\) and the duality isomorphism
\[
  \mathbb{R}f_*\mathbb{R}\underline{\Hom}(\mathbb{L}f^*E,\mathbb{L}f^*F)
  \simeq
  \mathbb{R}\underline{\Hom}(E,\mathbb{R}f_*\mathbb{L}f^*F).
\]
To see that the arrow of (7.1.2) is an isomorphism for \(E\) perfect, we reduce trivially, by localization and dévissage, to the case \(E=A\), in which it is a tautology.

\sgaStableTarget{sga6.e01.prop.7-2}\textbf{Proposition 7.2.} \uline{There exists} (cf.\ SGA 5 I 1.6), \uline{for} \(E\in\Ob\Dpx(S)\) \uline{and} \(F\in\Ob\Dpx^+(S)\)\uline{, a canonical functorial arrow}
\[
  E\longrightarrow\mathbb{R}\underline{\Hom}(\mathbb{R}\underline{\Hom}(E,F),F).
\]
\uline{It is an isomorphism when} \(E\in\Ob\Dpx(S)_{\mathrm{parf}}\) and \(F=A\).

